% =====================================================================
%  Hans Brøchner, Erindringer om Søren Kierkegaard.
%  Det nittende Aarhundrede. Maanedsskrift for Literatur og Kritik,
%  ed. Georg and Edvard Brandes, vol. 5 (October 1876 -- March 1877),
%  March 1877, printed pp. 337--374.  Published posthumously (Brøchner
%  d. 1875).  Dated by Brøchner in a line above the signature:
%  "(Nedskrevet i Tiden fra 27. Dec. 1871 til 10. Januar 1872.)"
%  TRANSCRIPTION (Danish) -- diplomatic, from the page images.
% ---------------------------------------------------------------------
%  ---- SOURCE -----------------------------------------------------------
%  Scan: ~/bibliotek/Brøchner, Hans/1877-erindringer-kierkegaard-nittende-aarhundrede.pdf
%        sha256 f3ee9adffbc41633731dfaa1efff8a6ae56b20ba9b9a83c2ffe8e54d0e3e0eb7
%        33037104 bytes, 486 PDF pages, body pages 313.056 x 537.4 pts.
%        Internet Archive item detnittendeaarhu1876unse (University of
%        Illinois copy, full view), downloaded 2026-09-22:
%        https://archive.org/details/detnittendeaarhu1876unse
%        The whole volume; this article is a small part of it.  Each page
%        is a 500 ppi 1-bit JBIG2 mask (which carries the type) over a
%        low-resolution colour image.  Embedded text layer (Internet
%        Archive OCR) -- the machine witness for ocrdiff.py; SCANS.tsv
%        records it USABLE.
%
%  ---- PAGE MAP ---------------------------------------------------------
%      PDF page = printed page + 10, uniform over printed 336--376.
%  Verified 2026-09-22 AT THE IMAGE: the folio was read off the rendered
%  leaf on every page PDF 346--386 (pp. 336, 338--374, 376); pp. 337 and
%  375 open a piece and print no folio.  No plate, advertisement or
%  unnumbered leaf falls inside the article.  Independently, the scan's
%  text layer gives offset 10 from 251 folios (paginate.measure_offset).
%  Sheet signatures at the page foot: 22 (p. 337), 22* (339), 23 (353),
%  23* (355), 24 (369), 24* (371) -- not transcribed.
%  Folio misprint: p. 356 (PDF 366) prints its folio "856" (the first figure
%  a closed 8, cf. the 8 of "358"); text and page map unaffected.
%
%  ---- WHAT SURROUNDS THE ARTICLE IN THE VOLUME -------------------------
%  Precedes (ends p. 336 = PDF 346): Edvard Brandes, review "Nye Samlinger
%    af Folkepoesi", signed "Edvard Brandes."; below it a rule and a
%    "Rettelse" (a correction to p. 159) -- the close of the preceding
%    monthly number.  P. 337 carries no issue heading; the top of the page
%    is blank and the article opens under a sunk head.
%  Follows (begins p. 375 = PDF 385): "Om nyere spansk Literatur.", by
%    S. Schandorph (so the running head of p. 376).
%  Neither is transcribed.
%
%  ---- LAYOUT AND TYPE --------------------------------------------------
%  Body in ANTIQUA (roman), about 35 lines to the page -- not Fraktur.
%  Head: "Erindringer om Søren Kierkegaard." in bold sans-serif, over a
%    short rule; no byline, no editorial headnote from the Brandeses, no
%    footnote at the head.
%  The article is divided into numbered sections, the number set centred
%    on its own line ("1.", "2.", ...); set here with \secnum{N.}.
%  Running heads "H. Brøchner: Erindringer om Søren Kierkegaard." -- not
%    transcribed.
%  Close: the dateline, centred, upright roman in smaller type, then the
%    signature "H. Brøchner." set right in bold condensed type.
%
%  ---- CONVENTIONS ------------------------------------------------------
%  Diplomatic.  Orthography, punctuation and capitalisation of 1877 as
%    printed; nothing modernised, nothing harmonised.  This journal prints a
%    comparatively modern spelling (gøre, Kristendommen, Filosofi, Øjeblik)
%    beside aa and capitalised nouns; the article has no å.
%  Greek is set in Greek (Διαψάλματα p. 348, μεσότης p. 361).
%  Omission points and withheld names are set as printed (spaced points,
%    asterisks, leader lines on p. 352).
%  Emphasis: LETTERSPACING -> \emph{}.  Judged by measured letter advance
%    within the line (at 400 dpi letterspaced runs show gaps of ~10--11 px
%    between letters against 2--6 px in plain words on the same line;
%    calibrated on "hele", p. 365), never by impression.  Italic or other
%    type set against the roman -> \textit{}; bold -> \textbf{}.
%  Quotation marks „...“ as printed.  Book titles are set in quotation
%    marks, not italic.
%  Footnotes: *) set with \footnote at the mark.
%  A word broken across a printed line is joined (Cau\-%); a hyphen printed
%    at a line end as part of the word is kept, with a comment.
%  A word divided across a PAGE turn carries \opage{N} at the point of
%    division: "inden\opage{30}for".
%  Printer's errors are transcribed AS PRINTED and logged in a % comment at
%    the site, on its own line before the line concerned; nothing is
%    silently corrected.
%
%  ---- ACCURACY  [2026-09-23] ------------------------------------------
%  Word-accuracy: every batch diffed against the embedded layer before
%  splicing (ocrdiff.py --frag, offset 10). 4 batches (pp. 337-346,
%  347-356, 357-366, 367-374), 63 candidates, 0 transcription errors
%  corrected, 63 the witness's fault, 0 unresolved.  See
%  TRANSCRIPTION-PLAYBOOK.md §3 step 4.  Because the transcribers took their
%  words from that layer, the diff cannot see a layer error they copied;
%  the readings below exist for that class.
%
%  Method.  (1) Transcription: one agent per batch, words from the text
%  layer corrected at the image (400 dpi; 500 dpi for any glyph in doubt),
%  structure, emphasis (measured) and punctuation from the image.
%  (2) SECOND READING of ALL 38 pages, word by word and mark by mark at the
%  image, by four readers who had transcribed nothing and were forbidden the
%  text layer: 0 word-level discrepancies; 9 rows, all faint marks or
%  punctuation (5 of them undecidable on the PDF's 1-bit images).
%  (3) THIRD READING by a reader who had done neither: settled the disputes,
%  and re-read pp. 343, 360, 369 in full as a control (seeded draw,
%  20260922): 0 discrepancies of any kind.
%  (4) The PDF's page images are 1-bit masks thresholded from Internet
%  Archive's captures, and faint ink is lost in them.  The disputed marks
%  were therefore settled finally on IA's processed full-tone JP2s of the
%  same pages (pp. 337, 342, 349, 351, 358, 364, 368, 374; fetched
%  2026-09-23 from the jp2.zip of the same item; scandata.xml confirms
%  PDF page = leaf number and IA's own page numbering = PDF - 10).  That
%  overturned the second and third readers once: "Halvtredserne," (p. 358)
%  is a comma whose tail the 1-bit image had lost.
%
%  Result: 0 word errors found in 38 pages (95% upper bound on the rate
%  ~0.08 per page).  Corrected after reading: p. 337 "Grupper:" -> "Grupper."
%  (a colon is set off by a thin space in this fount; this mark is not);
%  p. 358 "Universitetet" -> "Universitetet," (a sort's set width, head
%  unprinted); p. 368 "ham.." -> "ham." (the second "point" is a speck
%  taking no set width; the logged printer's error was withdrawn);
%  p. 374 "Gòtz" -> "Götz" (one diaeresis dot printed).  Apostrophes keyed
%  as straight ' on pp. 347-372 were normalised to ’, the glyph the page
%  prints throughout (23 sites; encoding only).  No silently corrected
%  printer's error was found, so none had to be restored.
%
%  Printer's errors logged at the site (as printed): p. 337 "K.‘s" turned
%  comma; p. 340 a second section "3." (no 4); p. 341 "erindrer;", "eiler";
%  p. 342 "i hvilke" (probable; sense wants hvilken); pp. 348, 350 opening „
%  missing; p. 352 "8krift", "af han" (at), „De skal ... never closed;
%  p. 364 "Iudtrykket" (turned n); p. 365 "sammen, Engang"; p. 366
%  "Forboldet"; p. 374 „ printed as two commas, „Godt! ... never closed.
%  Quote balance 0 („72 / “72) only because those defects cancel.
%
%  OPEN: p. 358 "Universitetet," -- comma or point undecidable even on the
%  JP2 (comma read; see the site).  p. 374 "Götz" -- ö or ò undecidable on
%  the JP2 (ö read).  p. 337 "K.‘s" was not re-examined on the JP2.
%  The speck-sized marks at p. 349 "er", p. 351 "det", "begges", "endnu"
%  (which the text layer reads as ér, dét, and commas) are settled as
%  specks on the JP2 and noted at each site.
%
%  Compile: the §5 recipe (cloud variant: lmodern absent there, so
%  libertinus, libertinust1math, babel and microtype's expansion stripped;
%  textalpha kept, so the Greek is exercised): 0 errors, 0 missing
%  characters, 38 distinct marginal page numbers.  A clean compile is a
%  typesetting result, not evidence of accuracy.
% =====================================================================
\documentclass[12pt,a4paper,openany]{book}
\usepackage[utf8]{inputenc}
\usepackage[T1]{fontenc}
\usepackage{textalpha}   % Greek: Διαψάλματα (p. 348), μεσότης (p. 361)
\usepackage{libertinus}
\usepackage{libertinust1math}
\usepackage{graphicx}    % for \reflectbox (the old Danish ɔ: id-est mark)
\DeclareUnicodeCharacter{0254}{\reflectbox{c}}  % ɔ (U+0254) = reversed c
\usepackage[danish]{babel}
\usepackage[protrusion=true,expansion=true]{microtype}
\usepackage{setspace}
\usepackage[margin=1.2in]{geometry}
\usepackage{fancyhdr}
\setlength{\headheight}{28pt}
\addtolength{\topmargin}{-13.5pt}

% Footnotes as the 1877 printing sets them: *).
\renewcommand{\thefootnote}{*)}

% The article's numbered sections: the number alone, centred on its line.
\newcommand{\secnum}[1]{\par\begin{center}#1\end{center}\par\nobreak}

\usepackage{hyperref}
\hypersetup{
  pdftitle={Erindringer om Søren Kierkegaard},
  pdfauthor={Hans Brøchner},
  hidelinks
}

\pagestyle{fancy}
\fancyhf{}
\fancyhead[LE,RO]{\thepage}
\fancyhead[RE]{\textit{H. Brøchner: Erindringer om Søren Kierkegaard}}
\fancyhead[LO]{\textit{Det nittende Aarhundrede, Marts 1877}}

\setstretch{1.3}

\title{\textbf{Erindringer om Søren Kierkegaard}}
\author{Hans Brøchner}
\date{Det nittende Aarhundrede. Maanedsskrift for Literatur og Kritik,
  5.\ Bind, Marts 1877, pp.~337--374}

\usepackage{marginnote}
\newcommand{\opage}[1]{\marginnote{\footnotesize\textit{[#1]}}}
\newcommand{\apage}[1]{\marginnote{\footnotesize\textit{[#1]}}}

\begin{document}
\maketitle
\thispagestyle{empty}
\clearpage

% --- p. 337 ---
% p. 337 prints no folio (opening page of the article); sheet signature „22“ at foot not transcribed.
\begin{center}\textbf{\large \opage{337}Erindringer om Søren Kierkegaard.}\\[4pt]\rule{2cm}{0.4pt}\end{center}

Jeg har villet optegne nogle af de Erindringer, jeg
fra en lang Aarrække har om Søren Kierkegaard, og
som kunne have nogen Interesse ogsaa for Andre. De
ville maaske engang, naar S. K. finder en kyndig Biograf,
kunne bidrage til at gøre hans Billede i enkelte Retninger
mere udført og mere levende. Jeg optegner mine
Erindringer netop som Materiale til et Billede af K. og
derfor uden alt Forsøg paa, allerede nu at bearbejde
dem til en Helhed eller at give dem en Afrunding ved
% „Grupper.“: a point, not a colon, though the text layer reads a colon. Settled by the third reading and
% confirmed on IA's processed JP2 of the same page (full-tone, 500 ppi; the PDF's 1-bit mask loses faint ink): the point stands ~5--7 px after the r, as every point after a letter does, whereas every
% colon in the article is set off by a thin space (12--23 px); the faint mark above it is a speck (2x3 px).
at samle dem i Grupper. De blive givne saaledes som
% emphasis: „enkelte“ letterspaced (letter gaps 10--13 px at 400 dpi; plain words on the line 2--6 px)
de ere bevarede, efter det første Indtryk, som \emph{enkelte}
Minder i en trofast Hukommelse.

\secnum{1.}

% „K.‘s“: the apostrophe is printed as a turned comma (‘), not the ’ used elsewhere (e.g. „K.’s“ p. 340); transcribed as printed.
Mit første Møde med S. K. falder i den allerførste
Tid efter at jeg var bleven Student (1836). Jeg saa’
ham da i min gamle Onkel, Grosserer M. A. Kierkegaards
Hus, i et Selskab, der blev gjort i Anledning af nuværende
Biskop Kierkegaards Forlovelse (eller Bryllup?) med hans
første Kone. Jeg saa’ der S. K. uden at vide, hvem
han var; der var kun blevet sagt mig, at han var Dr.
K.‘s Broder. Han talte kun meget lidt den Aften; forholdt
sig vistnok væsentlig iagttagende. Jeg fik kun et
% --- p. 338 ---
\opage{338}bestemt Indtryk af hans Ydre, over hvilket jeg nærmest
morede mig. Han var dengang 23 Aar gammel, havde
noget meget Uregelmæssigt i sin hele Skikkelse og en
mærkelig Koiffure. Hans Haar stod i en purret Kam
næsten et Kvarter op over hans Pande, og han saae
besynderlig forvildet ud derved. Jeg fik, uden at jeg
véd hvorledes, den Forestilling, at han var Handelskommis
--- maaske fordi Familien var en Købmandsfamilie, ---
og jeg supplerede det under Indtrykket af hans aparte
Ydre, straks til, at han maatte være i en Manufakturbutik.
Jeg har siden tidt let hjærtelig over min Skarpsindighed.

\secnum{2.}

Saa mindes jeg, i den nærmeste Tid, engang imellem
at have seet S. K. paa samme Sted. Jeg kom nu til at
tale med ham, og mærkede snart, at han hørte andensteds
hjemme end bag en Disk og med Alenen i Haanden.
Han morede sig blandt de meget jævne Omgivelser der
med at fremsætte Smaa-Paradokser. Saaledes erindrer
jeg, at han engang forbavsede sine Slægtninge ved, med
det alvorligste Ansigt af Verden at erklære, at han fandt,
at den gamle A. B. C., vi brugte i vor Barndom, var en
af de allerinteressanteste Bøger, der gaves, og at han
meget jævnlig endnu læste deri, og havde meget Udbytte
deraf. --- En Aften spillede han Boston med. Det spilledes
der med Variationer, af hvilke en bestod deri, at
der til 3 af de Spillende gaves 17 Kort til hver, til den
fjærde derimod kun et, og at saa de 3 andre hver afgav
4 Kort til den sidste. Der i Huset, hvor man tog Kortspillet
med andægtig Alvorlighed, betragtede man altid
den 4de Mands Situation i dette Spil som yderst sørgelig,
og S. K. forbavsede atter i høi Grad ved at udvikle, at
det var den allerfortræffeligste Situation, der kunde
tænkes, og at han kun ønskede bestandig at kunne være
den saaledes situerede fjerde Mand. Det vilde være det
Pikanteste af Alt.

% --- p. 339 ---
\secnum{3.}

\opage{339}I 1837 traf jeg af og til K. paa en Restavration.
Han boede dengang ikke længere hos Faderen, men i et
Hus i Løvstræde, vistnok det, hvor nu (1871) Reitzels
Boglade er. Sin Aftensmad tog han gærne paa en Restavration,
og jeg kan huske, at jeg, naar jeg engang
imellem gjorde mig tilgode med „en halv Boeuf“ forbavsedes
over den Luksus, S. K. udfoldede, naar han spiste
til Aften med en halv Flaske Vin o. desl. Vi talte nu
hyppigere sammen, og han viste sig meget venlig imod
mig. En Aften spurgte han mig, hvad jeg læste af
æstetiske Skrifter. Jeg kom i den Anledning til at
berøre, at jeg kun havde Adgang til Lidet, og at jeg
var meget ukendt paa mange Gebeter. Han spurgte
mig, om jeg kendte de tyske Romantikeres Skrifter,
hvilket jeg maatte benægte. Han opfordrede mig saa
til at følge hjem med sig og laante mig en Bog af
Eichendorff: Dichter und ihre Gesellen, --- en Bog, jeg
siden købte paa Avktionen over hans Bøger til en Erindring
om dette Møde med ham. Jeg mindes saa, at da
jeg efter en 14 Dages Tid bragte ham den tilbage, og
netop vilde gøre en Undskyldning fordi jeg, der dengang
studerede Teologi med Lidenskab, havde beholdt den
saalænge, blev jeg konfunderet ved, at han modtog mig
% emphasis: „allerede“ letterspaced (letter gaps 8--14 px at 400 dpi)
med det Spørgsmaal, om jeg \emph{allerede} havde læst den.
Han har øjensynlig, med sit skarpe Blik, kunnet se, at
det modsatte Ord svævede mig paa Tungen, og han har
moret sig med, paa en godmodig Maade at gøre mig
forlegen; hvilket dengang var overmaade let. Jeg kom
saa let til at rødme, og det Syn har altid tiltalt ham
hos „det unge Menneske“. --- Fra mit Besøg hos K., da
han laante mig Bogen, mindes jeg to Ting. Den ene er
min Forbavselse over hans store Bibliotek, der ganske
imponerede mig. Det andet var en Særhed hos K. Han
gik atter ud, da vi havde hentet Bogen og pustede, da
vi gik bort, et Vokslys ud, han havde tændt. Han for%
% --- p. 340 ---
\opage{340}klarede mig saa, at han altid gjorde det med en vis
Forsigtighed og paa nogen Afstand, fordi han havde
reflekteret sig ind i den Forestilling, at Osen af Lyset
var farlig at indaande, og at den kunde skade hans
Bryst.

% PRINTER'S ERROR: this section is numbered „3.“ as printed; it should be 4 (3 stands on p. 339, 5 follows below; no section 4 is printed).
\secnum{3.}

I 1838 døde S. K.’s Fader. Ham har jeg kun set
én Gang, hos min gamle Onkel Kierkegaard. Det var
en betydelig Personlighed, hvis Ydre indprægede sig levende
i min Erindring. Han gik med min Onkel op og ned
ad Gulvet i Onkels Spisestue, i Samtale med ham. Han
gik noget bøjet; hans Holdning svarede til hans Ansigts
tænksomme, alvorlige Udtryk; det graa Haar var strøget
tilbage bag Øret, hvorved Ansigtsformen traadte skarpt
frem. --- Den Ærefrygt og Kærlighed, S. K. nærede for
sin Fader, har fundet sit levende Udtryk i hans Skrifter.
--- Den Gamle var en meget punktlig, i Pengesager nøjagtig
Mand, sparsommelig til dagligdags, men ved given
Lejlighed rundhaandet. S. K. fortalte mig engang, hvorledes
han som ung Student havde gjort en Rejse i Nordsjælland,
og hvorledes den Gamle havde udstyret ham
rundeligt med Rejsepenge, og saa endda overrasket ham
med, til en af de sidste Rejsestationer at sende ham et
Pengebrev med 50 Rdlr. Om Faderens Punktlighed brugte
S. K. engang i en Samtale med Mag. Adler det Udtryk:
„min Fader er født i Terminen“. --- Da S. K. forlod
Faderens Hus, gav den Gamle ham aarlig en Sum, der
efter den Tids Forhold var meget rigelig; tager jeg ikke
meget feil, var det 800 Rdlr. om Aaret.

\secnum{5.}

Min afdøde Søster Hansine var i 1839 som ung Pige
her i Byen. Engang da jeg kom til at nævne S. K.,
fortalte hun mig, at hun havde set ham, da hun som
ukonfirmeret var her i Byen i Besøg hos gamle Onkel
% --- p. 341 ---
\opage{341}Kierkegaard. Vore Kusiner havde en Dag, da S. K.
skulde komme der, givet hende det Raad, ikke at indlade
sig med ham, da han var „en forskrækkelig forkælet
og uartig Dreng, der altid hang i sin Moders Skørter.“
Min Søster var dengang 13 Aar, S. K. 15 Aar gammel.
% „nogle“: a speck stands over the e (reads like ė); not a letter.
Hun læste sildigere med stor Glæde nogle af hans opbyggelige
Skrifter, navnlig „Kærlighedens Gerninger“, og
kunde saa undertiden, med det for hende ejendommelige
Smil, minde mig om Kusinernes Formaning til hende.

\secnum{6.}

Kort efter at S. K. i 1840 var bleven teologisk
Kandidat, talte jeg med ham. Han fortalte mig, at
hans Fader bestandig havde ønsket, at han skulde tage
teologisk Examen, og at de meget jævnligt havde diskuteret
den Sag. „Saalænge Fader levede, kunde jeg
imidlertid forsvare min Teori, at jeg ikke burde tage
den, men da han var død, og jeg ogsaa maatte overtage
% „partes“ is set in italic.
hans \textit{partes} i Debatten, saa kunde jeg ikke længere staa
mig, og saa maatte jeg beslutte mig til at læse til Eksamen.“
Dette havde S. K. saa gjort med stor Energi. Han havde
taget „Sin Brøchner“ til Manuduktør, havde arbeidet sig
gennem de tørreste Discipliner, gjort Ekcerpter af Kirkehistorien,
lært „Paverækker“ udenad o. desl. Hans Manuduktør
var meget glad over ham. Kort efter at S. K.
havde taget Eksamen, kom Peter Stilling (senere Dr. phil.)
til Sin Brøchner og vilde manuduceres. Han ytrede, at
han mente at kunne blive færdig paa halvandet Aar;
„længere havde S. K. jo ikke læst.“ „Ja vist“ sagde
gamle Br., der ikke ekcellerede i Høflighed, „bild Dem
ikke det ind! med S. K. var det en anden Sag, han
% PRINTER'S ERROR: „erindrer;“ -- semicolon printed where a comma is wanted; transcribed as printed.
kunde Alting!“ --- Jeg erindrer; at jeg i 1839 har gaaet
paa teologiske Skriveøvelser sammen med S. K. hos Clausen.
% PRINTER'S ERROR: „eiler“ for „eller“ -- the second letter is a short dotted i, not l (checked at 500 dpi); transcribed as printed.
Han deltog imidlertid kun deri de første 2 eiler 3 Timer.
Der fortaltes, at han engang flere Aar tidligere havde
deltaget i Clausens Skriveøvelser og engang, istedenfor
% --- p. 342 ---
\opage{342}at besvare Opgaven, havde analyseret den og paavist, at
der ikke var Mening at finde deri. Derover var Cl. bleven
fortørnet, og S. K. havde ophørt at gaa paa Skriveøvelser.
--- Ved den skriftlige teol. Eksamen stod S. K. som Nr.
4 ved Censuren. Foran ham stode M. Wad, Warburg
og Chr. F. Christens. Censorene erklærede imidlertid
til disse, at K.’s Afhandlinger vidnede om en langt større
Modenhed og Udvikling af Tanker end nogen af de Andres;
men disses indeholdt et rigere Maal af positivt teologisk Stof.

\secnum{7.}

% PRINTER'S ERROR (probable): „i hvilke“ with the singular „Samtale“ -- „hvilken“ expected; transcribed as printed, no n in the print.
I den samme Samtale, i hvilke S. K. fortalte mig,
hvorledes han var kommen til at tage teol. Examen,
omtalte han ogsaa sin Faders forunderlige Ro og Objektivitet.
Saaledes havde den Gamle engang sagt til S. K.:
„det var egentlig en Lykke for Dig, om jeg var død;
saa kunde Du maaske endnu blive til Noget; det bliver
Du ikke, saalænge jeg lever.“ --- (Et andet Træk af hans
Fader blev mig fortalt af min afdøde Fætter Peter Kierkegaard.
% „S. K.’s“: the point after S has a faint mark above it (reads like „S:“); taken as a speck.
Da S. K.’s Moder døde, var hans Fader dybt
bedrøvet. Ved Ligbegængelsen kom Biskop Mynster hen
til ham og udtalte bevæget sin Deltagelse. Den Gamle
hørte, uden at røbe sin Sorg, paa Mynster, hvem han
skattede højt, og da Mynster tav, sagde han: „Skal vi
saa ikke, Deres Højærværdighed, gaa ind i Værelset der
og drikke et Glas Vin.“ Mynster, der kendte ham, misforstod
ikke denne Ytring.)

\secnum{8.}

Medens S. K. var i Berlin 1841---42, var det, at jeg
indstillede mig til teologisk Eksamen, hvortil der paa Grund
af mine uortodokse Anskuelser blev nægtet mig Adgang.
% „med“: the first stroke of the m is broken (reads like „ıned“).
Jeg korresponderede dengang med Chr. F. Christens, der
ogsaa var i Berlin, og gennem ham sendte og modtog
jeg Hilsener, ogsaa til og fra Kierkegaard. Kort efter
hans Hjemkomst traf jeg ham paa Gaden. Han var
% --- p. 343 ---
\opage{343}meget venlig imod mig, talte om min Stilling og mine
Studier, og sagde mig, hvad der glædede mig, at han i
Berlin havde længtes efter flere Ting herhjemme, deriblandt
ogsaa efter at se mig. Sligt kunde han sige med
et ejendommeligt vindende Udtryk. Hans Smil og hans
Blik var ubeskrivelig udtryksfuldt. Han havde en egen
Maade, paa Afstand at hilse med et Blik. Det var kun
en lille Bevægelse af Øiet, og dog udtrykte det saa Meget.
Der kunde være noget uendeligt Mildt og Kærligt i hans
Blik, men ogsaa noget Æggende og Tirrende. Han kunde
ved et Blik paa en Forbigaaende uimodstaaelig „sætte
sig i Rapport“ til ham, som han udtrykte det. Den, der
mødte dette Blik, blev enten tiltrukken eller frastødt,
gjort forlegen, usikker eller tirret. Jeg har gaaet en
hel Gade igennem med ham, medens han udviklede mig,
hvorledes man, ved saaledes at sætte sig i Rapport til
de Forbigaaende, kunde gøre psykologiske Studier, og
medens han udviklede Teorien, realiserede han den i
Praksis omtrent med hvert Menneske vi mødte. Der
var ikke en, paa hvem hans Blik ikke gjorde et synligt
Indtryk. Ved samme Lejlighed forbavsede han mig ved
den Lethed, med hvilken han knyttede en Samtale med
mangfoldige Mennesker, i nogle faa Repliker optog en
tidligere Samtale og førte den et Skridt videre til et
Punkt, hvor den ved Lejlighed atter kunde optages.
Anledningen til disse Eksperimenter var, at jeg en Dag
havde gaaet foran ham, fordybet i Tanker, saa at jeg
ikke havde hørt, at han kaldte paa mig, og ikke mærket
at han pikkede mig paa Skuldren. Da jeg endelig mærkede
ham, sagde han, at man gjorde Uret i, saaledes at synke
ind i sig selv og ikke at gøre de Iagttagelser, man i
saa rigt Maal kunde gøre. For saa at vise mig Metoden,
trak han mig med sig op og ned ad flere Gader, og forbavsede
mig med sit psykologiske Eksperimenteretalent.
Han var altid interessant at følges med, men i én Henseende
havde det sin Besværlighed at følges med ham.
% --- p. 344 ---
\opage{344}Paa Grund af det Uregelmæssige i hans Bevægelser, der
vistnok hang sammen med hans Skævhed, kunde man
% „Linje naar“ and, below, „endnu“: type damaged (broken n in each), words legible.
aldrig holde den rette Linje naar man fulgtes med ham;
man blev altid sukcessivt trængt ind efter mod Husene
og Kælderhalsene, eller ud efter mod Rendestenen. Naar
han saa tilmed gestikulerede med Armen og sin Spanskrørsstok,
blev det endnu mere en Marche med Forhindringer.
Man maatte af og til benytte en Lejlighed til
at komme om paa den anden Side af ham forat vinde
den fornødne Plads.

\secnum{9.}

Den Vinter S. K. tilbragte i Berlin, traf han der
sammen med mange Danske, af hvilke ikke faa vare
komne dertil for at høre Schelling. Blandt dem var
min afdøde Ven Christian Fenger Christens. Ham omtalte
S. K. med stor Anerkendelse; han udtalte senere
for mig, at Christens var den dygtigste af alle de Danske
der den Vinter havde været i Berlin, og blandt dem
vare dog saadanne Mænd som Lic. jur. (senere Minister)
A. F. Krieger, og Avditør (nu Departementschef i Kultusministeriet)
Carl Weiss. Christens fortalte mig forskellige
Smaatræk om K., om hans Ubehjælpsomhed naar han
skulde tale Tysk, navnlig naar han skulde tale om Ting
af det daglige Liv (saaledes var det ham den første
Aften i Berlin ikke muligt at faa det udtrykt, at han
vilde have en Lysestage); --- om hvorledes hans Vært
trak ham op og eftersom Optrækkerierne tiltog, til Gengæld
lod ham avancere fra Magister til Doktor og saa
til Professor; --- Noget, Kierkegaard selv morede sig
over, da han nok anede Grunden, --- om hans Forlegenhed,
da Weiss, der satte Pris paa en elegant Middag, en Dag
havde ført ham ind paa en Restavration af første Rang,
hvor et Antal af elegantklædte Herrer i sorte Kjoler,
med hvide Halstørklæder og Halsbind stode i Hjørnerne
af Salen, og efterat K. ærbødigt havde hilst paa dem
% --- p. 345 ---
% The asterisks below stand in the print for a suppressed name.
\opage{345}og med W. sat sig ved et Bord, kom springende og viste
sig at være Opvartere o. s. v. --- Christens fortalte mig
ogsaa, hvorledes Kierkegaard i Berlin morede sig med
at faa Lic. theol. * paa Glatis. * havde allerede været
udenlands i længere Tid, havde hørt Forelæsninger i
Strasburg og andensteds og var meget selvbevidst i
Anledning af al den Visdom, han havde indsamlet. S. K.
undlod sjelden, naar de traf sammen, at bede ham om
at meddele sig de væsentligste videnskabelige Resultater
af hans Udenlandsrejse. Naar saa * formulerede en eller
anden Sætning som vundet Udbytte, forstod K. altid,
enten at vise ham, at det var noget meget Gammelt,
eller ved at gøre smaa Spørgsmaal om „Noget, der ikke
var ham ganske klart“ i det Fremsatte, at bringe den
stakkels Licentiat, der kun var lidet begavet som Tænker,
ind i en saadan Konfusion, at han hverken vidste ud
eller ind. Naar K. saa havde faaet ham paa det Punkt,
kom han pludselig i Tanker om, at han havde meget
travlt og nødvendigvis skulde et eller andet Sted hen.
Saa lod han Licentiaten staa midt i Konfusionen til stor
Moro for de Andre. --- I Enten-Eller hos K. skildres
denne Metode, netop anvendt paa en Licentiat. Mon
han ikke skulde have havt sit Berlinske Bekendtskab i
Erindring da han skrev hine Sider?

\secnum{10.}

Ikke længe efter K.’s Hjemkomst talte han engang
med mig om mine Studier. Jeg fortalte ham hvad jeg
paa den Tid læste af græske Forfattere og berørte den
Interesse, der mere og mere vaagnede hos mig for den
græske Poesi og Filosofi. Han opmuntrede mig til at
fortsætte disse Studier og talte med Varme om Betydningen
af det Græske for vor Tid. Han brugte det Udtryk:
„Grækenland burde bestandig have en Repræsentant
hos vor Tid, ligesom den ene Magt har en Repræsentant
hos den anden til at gøre sine Interesser gæl%
% --- p. 346 ---
\opage{346}dende.“ --- K. havde selv gode Kundskaber i Græsk og
havde læst ikke Lidet. Hvor fint han opfattede den
græske Aand, bære hans Skrifter Vidnesbyrd om.

\secnum{11.}

I 1842 udkom et lidet pseudonymt Skrift: „Johan
Ludvig Heiberg efter Døden.“ Jeg kendte til Skriftet
før det udkom, og omtalte det engang for S. K. Han
var ikke fornøjet med, at Heiberg blev gjort til Genstand
for en saadan Spøg og udtalte sig med megen
Varme om hans Betydning som Æstetiker og navnlig
% footnote *) on p. 346
som sin Generations Læremester i Æstetik i vort Fædreland\footnote{K. stillede H. over alle de samtidige Æstetikere i Tyskland (NB. Vischer var dengang endnu ikke optraadt uden med mindre Ting.)}. Senere efterat Enten-Eller var udkommen, og
Heiberg havde skrevet sin bekendte Anmeldelse af den,
talte S. K. engang til mig om ham og lagde sin Misfornøjelse
over H.’s Optræden tydelig for Dagen. Han
anerkendte H.’s Betydning som Æstetiker, men akcentuerede
nu ogsaa stærkt hans Begrænsning. „Jeg vilde
kunne drage en hel Række af æstetiske Problemer frem,
om hvilke H. ikke har en Anelse.“ Han sigtede hermed
vistnok til de Problemer, der ere komne frem i „Stadier
paa Livets Vei“ og „Afsluttende Efterskrift“.

\secnum{12.}

Jeg mødte i de Aar hyppig S. K., naar jeg om Aftenen
gik en Spadseretur. I Regelen var Frederiksberg
Have Maalet for min Vandring. Saalangt gik ogsaa K.,
men kun til Indgangen af Haven, hvor de smaa Blomsterpartier
ere paa hver Side af den dengang smalle Vej,
der fører fra Porten til den første frie Plads. Han indaandede
der i nogle Øjeblikke Duften af Blomsterne, og
% --- p. 347 ---
\opage{347}tog saa Erindringen om dette „Øjeblik“ med sig. Saaledes
holdt han af at afslutte og begrænse: hans Vandring havde et
bestemt Maal, men dette Maal blev ligesom kun berørt; der blev
ikke dvælet ved det og der toges ligesom kun et Nydelsens Motiv,
der kunde ideelt bearbejdes.

Til en af disse Aftenvandringer knytter sig en lille komisk
Erindring. Da vi gik tilbage fra Haven til Alleen og skulde gaa
imellem de Afvisere, der staa ved Enden af Vejen imellem Træerne,
der er bestemt for Fodgængere, vare Afviserne nylig malede. K.
lagde ikke Mærke dertil, men støttede Haanden paa dem, idet han
gik imellem dem. Han mærkede saa, at hans Haand var fuld af
Maling, og gjorde nu, medens han fortsatte sin Samtale, en Række
af uheldige Forsøg paa, med den uundværlige Stok under Armen, at
vaske Malingen af ved at gnide de af Duggen fugtige nedhængende
Blade af Træerne mellem Hænderne. Hans Anstrængelser fremkaldte
mangt et Smil paa de Forbigaaendes Ansigter.

\secnum{13.}

Paa den Tid saa’ jeg ogsaa af og til S. K. til Hest. Han havde
lært at ride for at kunne tage sig denne Motion og for at kunne
gøre Smaature uafhængig af Kuske o. desl. Han gjorde ingen
synderlig god Figur paa Hesten. Man saa’ paa hans Holdning, at
han ikke tiltroede sig nogen stor Evne til at styre den, hvis den
skulde faa isinde at gøre sig rebelsk. Han sad stiv paa Hesten,
og man fik Indtrykket af, at han stadig mindedes Ridelærerens
Instrukser. Nogen videre Frihed til paa Hesteryggen at forfølge
sine Tanker og sine Fantasier kan han neppe have havt. Han opgav
ogsaa snart denne Idræt og kørte saa hellere, naar han vilde
besøge sine Yndlingssteder i Skovene i Københavns Omegn. I de
Aar, da hans Forfattervirksomhed var intensivest, vare disse Ture
et af de Midler, han anvendte for at holde
% --- p. 348 ---
\opage{348}sig frisk, og for at bringe sig i den Stemning,
Produktionen udfordrede.

\secnum{14.}

Et enkelt Udtryk af K. fra hin Tid anfører jeg som karakteristisk.
Han havde en Tjener, der var meget paalidelig og til hvem K.
betroede Udførelsen af alleslags praktiske Ting. Naar han saaledes
skulde flytte, kørte han om Morgenen ud og om Aftenen drog han
hjem i sin nye Lejlighed, hvor Alt stod i fuldkommen Orden --- selv
Biblioteket var ordnet --- ved Tjenerens Hjælp. Om denne Tjener
sagde K. engang, da han havde omtalt hans praktiske Talent: „han
er i Virkeligheden \emph{mit Legeme}.“

\secnum{15.}

I den første Tid efter at Enten-Eller var udkommen, havde jeg
hyppige Samtaler med ham om Bogen. Jeg havde straks ved
Gennemlæsningen af Διαψάλματα genkendt den anonyme Forfatter, jeg
havde straks fundet Ord, der vare faldne i tidligere Samtaler med
K. Jeg forsøgte naturligvis aldrig, i mine Samtaler med K. at faa
ham til at udtale sig som Forfatter; men hvor Noget var mig
dunkelt, talte jeg med ham derom som med den Ældre, der havde den
rigere Indsigt. Saaledes fik jeg mangen en nærmere Udvikling og
han udtalte sig med en Uforbeholdenhed om Bogen, som han neppe
vilde have vist, hvis jeg havde været indiskret. Et Punkt, ved
hvilket han flere Gange dvælede, var det poetiske Element i 1ste
Del. Han udviklede en Dag med meget Liv, hvorledes der paa mange
Steder var givet Motivet til et Digt, men et Motiv, der vist med
Forsæt ikke var udført. Saaledes fremhævede han som Eksempel
Skildringen af Kvindens Skønhed og sagde:
% PRINTER: no opening „ before "i hver af disse Sætninger"; the closing “ after "Sonet." is printed. As printed.
i hver af disse Sætninger: den lille Haand, den nydelige Fod, o.
s. v. ligger der et Motiv til en Sonet.“ --- Han antydede af
% --- p. 349 ---
\opage{349}og til ligesom Forlængelseslinjerne af de Problemer,
Enten-Eller behandlede; men flere af hans Antydninger bleve mig
først ret forstaaelige, da de efterfølgende Skrifter udkom. --- Han
var ikke uimodtagelig for anerkendende Ytringer om Bogens
Betydning. Jeg ytrede en Dag, at ingen Bog, siden jeg for første
Gang læste Hegels Logik, havde sat mine Tanker i Bevægelse
% line-end hyphen of the compound Enten-Eller kept
saaledes som Enten-Eller. Denne Ytring tiltalte ham øjensynlig. Vi
stode netop ved Døren til hans Bolig (paa Nørregade) og vare ifærd
med at gaa hver til Sit; han gav mig, da vi skiltes ad, venligt
Haanden med det Smil, jeg kendte saa godt.

\secnum{16.}

Da jeg havde udgivet mit Skrift mod Martensen, („Nogle
Bemærkninger om Daaben“) talte K. med mig derom. Han var utilfreds
med en Anmeldelse i Berlingske Avis, der havde forsøgt at anslaa
en overlegen Tone, og ved en tiltagen Myndighed ligesom vilde
trykke den unge Forfatter ned. Jeg sagde ham, at jeg egentlig kun
havde moret mig over Anmeldelsen, da dens forcerede Kristelighed
fik et lidt komisk Anstrøg, naar man huskede paa, at den anonyme
Artikel stod i et Blad, der redigeredes af en Jøde. --- I mit lille
Skrift fremhævede K. en enkelt Ytring, som han vilde have ønsket
udeladt; det var, tror jeg, hvor jeg med nogen Spot over
Martensens Forsøg paa at hævde det nyfødte Barns Frihed, siger, at
den eneste Frihedsakt, der hos det nyfødte Barn kan være Tale om
% "er": a tiny dot over the e (the text layer reads "ér"). Settled on IA's processed JP2 of the same page (full-tone, 500 ppi; the PDF's 1-bit mask loses faint ink): a speck, not an acute (the acute of "véd" p. 351 is a 10x12 px wedge).
ved Daaben, er den, at Barnet skriger, naar det overøses med Vand.
Denne Ytring, mente K., saa’ kaad ud, og kunde let blive benyttet
af Vrangvillige.

\secnum{17.}

Ludvig Feuerbach omtalte K. hyppigere i vore Samtaler. Det er
efter mundtlige Ytringer af ham, at jeg i Problemet om Tro og
Viden (S. 125) har sagt at der ingen Tvivl kan være om, at det er
til ham, K. har
% --- p. 350 ---
\opage{350}sigtet ved et der anført Citat. Han anerkendte hos
Feuerbach den Klarhed og Skarphed, med hvilken han, netop som
Forarget, havde opfattet Kristendommen; men han fremhævede engang
i en Samtale, at han af Feuerbachs Naturbegejstring, særlig af den
Patos, med hvilken han i Wesen des Chr. taler om Vandets styrkende
Kraft, ogsaa i aandelig Henseende, havde faaet Indtrykket af ikke
blot en stærk Sanselighed hos F., men ogsaa af en Svækkethed ved
Sanselighed. Han føjede til:
% PRINTER: no opening „ before "dette er en Mening"; the closing “ after "Gisning." is printed. As printed.
dette er en Mening, jeg naturligvis aldrig vil udtale paa Tryk;
men jeg har ikke kunnet afværge dette Indtryk, og der er Noget i
Feuerbachs ejendommelige lidenskabelige Patos, og i hans Antipati
mod Kristendommen, der bliver mig tydeligere ved denne Gisning.“
--- Denne Ytring om F. bragte ham til at tale om Figurer fra
Oldtiden, hos hvem en raffineret Sanselighed traadte frem.
Saaledes gav han et med megen Fantasi udført Billede af Mæcenas,
hvorledes han, der „vistnok havde udtømt alle Roms Nydelser“, i
sin Søvnløshed søgte Bedøvelse ved Springvandets sagte Plasken i
hans Sovekammer og ved svage Toner af fjærn Musik. Saadanne
Billeder kunde han udmale fint; det var dog ikke den digteriske
Skildring, men den psykologiske Analyse, eller rettere Dissektion,
der interesserede ham.

\secnum{18.}

Paa en længere Spadseretur i de Aar fik jeg engang Lejlighed til
at iagttage, med hvilket Udtryk og hvilken Kraft K. kunde fremsige
Stykker af et Digterværk. Han talte om Aladdin, om den geniale
Dristighed i at ønske, der betegner ham, om Modsætningen mellem
denne Geniets kække Tillid og Refleksionens forsigtige Beregning.
Han citerede nu Aladdins Ord til Aanden, hvor han byder ham, til
Brudenatten at bygge det herlige Palads. Vi gik netop udenfor
Charlottenborg. K. standsede, da han fremsagde disse Linier, som
han ledsagede med
% --- p. 351 ---
\opage{351}en let Gestus med Armen. Der kom Farve i hans Kinder,
og hans Blik blev lyst og varmt, og hans Stemme var, skøndt
dæmpet, som en Herskers, der véd, at han har Magten til at byde.

\secnum{19.}

Engang talte vi om Hegels Fremstilling af Filosofiens Historie. K.
fremhævede det paa mange Punkter Subjektive i Opfattelse, og
sagde: „egentlig mangle Genier altid Evnen til en objektiv
Opfattelse af Andres Tanker; de finde overalt deres egne.“ ---
Ikke længe efter kom han i en Samtale med mig til at sige om sig
selv: „Jeg har aldrig havt denne Evne til at opfatte Andre
objektivt.“ --- Han erindrede i det Øjeblik neppe Ytringen i vor
tidligere Samtale; jeg mindedes den, og smilte ved mig selv, idet
jeg uddrog Konklusionen af hans Ytring om sig selv.

\secnum{20.}

Om Mag. Adler fortalte K. mig paa den Tid, da Adlers
Sindsforvirring begyndte, et Par kuriøse Træk. Adler kom en Dag
til K. med et Skrift, han havde udgivet og talte længere med ham
% "begges": a speck at the baseline (the text layer reads "begges,"). Settled on IA's processed JP2 of the same page (full-tone, 500 ppi; the PDF's 1-bit mask loses faint ink): no comma prints.
om begges religiøse Forfattervirksomhed. Adler lod K. forstaa, at
han betragtede ham som en Slags Johannes den Døber, i Forhold til
sig selv, der, som den, der havde modtaget en direkte Aabenbaring,
var den egentlige Messias. Jeg mindes
% "endnu": a speck after it (the text layer reads "endnu,"). Settled on IA's processed JP2 of the same page (full-tone, 500 ppi; the PDF's 1-bit mask loses faint ink): no comma prints.
endnu det Smil, med hvilket K. fortalte mig, at han havde svaret
Adler, at han var fuldkommen tilfreds med den Stilling, Adler
% "det": a tiny dot over the e (the text layer reads "dét"). Settled on IA's processed JP2 of the same page (full-tone, 500 ppi; the PDF's 1-bit mask loses faint ink): a speck, not an acute.
anviste ham: han fandt, at det var en meget respektabel Funktion
at være en Johannes den Døber og aspirerede ikke til at være
Messias. Under samme Besøg oplæste Adler en større Del af sit
Skrift for K. Noget deraf med sin sædvanlige Stemme, Andet med en
ejendommelig hviskende. K. tillod sig en Bemærkning om, at han
ikke kunde finde nogen ny Aaben%
% --- p. 352 ---
\opage{352}baring i Adlers
% PRINTER: "8krift" -- the initial S is printed as a figure 8 (wrong sort). As printed.
8krift, hvorpaa Adler sagde til ham: „Saa vil jeg igen komme til
Dig i Aften og læse hele Skriftet for Dig med \emph{denne} Stemme
(den hviskende), saa skal Du se, at det vil gaa op for Dig.“ K.
morede sig, da han fortalte mig Historien, meget over denne
Formening af Adler, at Variationen af Stemmen skulde gøre
Skriftet betydeligere. --- Af Adler hørte jeg engang en Ytring om
Kierkegaard i mit første Studenteraar. Han omtalte K.’s aandrige
Konversation, men fremsatte den Mening, at K. undertiden
forberedte sig paa sine Samtaler.

\secnum{21.}

K. havde som Yngre, og maaske endnu efter sin Disputats, tænkt
paa Muligheden af en Virksomhed ved Universitetet. Denne Tanke
traadte imidlertid snart tilbage. Engang fortalte han mig, at
Sibbern havde opfordret ham til at søge en Ansættelse som
filosofisk Docent.
% PRINTER: "af han saa" for "at han saa". As printed.
K. havde svaret, af han saa ialfald maatte betinge sig et Par
Aars Forberedelsestid. „Ih! hvor kan De tro, at man vil ansætte
Dem paa de Vilkaar?“ sagde saa Sibbern. „Ja, jeg kunde jo
ogsaa“, replicerede K., „gøre som Rasmus Nielsen, og lade mig
ansætte uden at være forberedt.“ --- Sibbern blev gnaven og sagde:
% PRINTER: the quotation opened at „De skal is never closed; "Nielsen! ---" ends it and the text runs on. As printed.
„De skal nu ogsaa altid saaledes hakke paa Nielsen! --- Det var
kort Tid efter, at Nielsen var bleven ansat ved Universitetet.
% The printed text ends here in a dash followed by leader dots filling the rest of the line and two further full lines of dots (an omission), then section 22.
---\,\dotfill\par
\noindent\dotfill\par
\noindent\dotfill

\secnum{22.}

K. har i forskellige af sine Skrifter udtalt sig om Korsarens
Angreb paa ham. Han har der væsentlig set Sagen fra et etisk
Synspunkt, og, med Hensyn til den Betydning, Korsarangrebene
havde for den almindelige Bevidsthed, har han vistnok i høj Grad
overvurderet
% --- p. 353 ---
\opage{353}den. K. havde ikke, i Forhold til et givet Fænomen, en
Virkelighedssans, om jeg maa bruge dette Udtryk, der kunde danne
Modvægten mod hans saa enormt udviklede Refleksion. Han kunde
ligesom reflektere en Bagatel op til at blive af verdenshistorisk
Betydning. Saaledes gik det ham vistnok med Korsaren. --- Til mig
har han i en Samtale omtalt hele Korsarens Retning og bedømt dens
Vittighedsmageri fra et rent æstetisk Synspunkt. Hans væsentlige
Anke mod den var, at den ødelagde Publikums Sans for det Komiske i
dets Idealitet ved altid at knytte Latteren til bestemte kendte
Personer, der nu engang vare slaaede fast for dens Læsere som
komiske Figurer, saa at Læseren lo af hvad der sagdes, ikke fordi
det virkelig indeholdt et komisk Gehalt, men fordi det, uden at
være virkelig komisk, sagdes om Personer, som man nu engang var
enig om at finde komiske og at le af.

\secnum{23.}

Noget Lignende, der hændte K. ved Korsarens Angreb, hændte ham
ogsaa med „Søren Kirk“ i Hostrups „Genboerne“. I sine efterladte
Dagbøger har K. hyppigere berørt det, at Hostrup har bragt ham paa
Scenen, og han har gjort det paa en Maade, der vidner om, at han
har følt sig afficeret derved. Jeg spillede, ved Stykkets første
Opførelse, denne Rolle, og jeg mindes, dengang at have talt med
K. derom uden at kunne mærke, at han lagde nogen Betydning deri;
men efterhaanden har han vistnok reflekteret sig ind i en anden
Stemning. Hostrup har i den lille Rolle taget Motiver fra S. K.’s
Skrifter, men han har ved Rollen øjensynlig nærmest tænkt paa den
Art Dialektik, der i de Aar grasserede blandt de unge Studerende,
% five spaced points as printed, standing for a name withheld
efterat ..... havde givet Impulsen til en meget overfladisk
Beskæftigelse med Filosofien. I denne Art af Dialektik var K. i
sine yngre Aar Mester --- naar han vilde være det. --- Min Samtale
% --- p. 354 ---
\opage{354}med K. om Rollen fandt Sted da jeg en Aften gik til en
Prøve paa Hofteatret. Paa Højbroplads mødte jeg K., der standsede
mig og spurgte, hvor jeg skulde hen. Jeg sagde ham, at jeg skulde
til Prøven paa Hostrups Stykke. „Ja, De skal jo spille mig!“
sagde saa K. Jeg svarede ham, at jeg just ikke kom til at spille
\emph{ham}, men en Dialektiker af vor hjemlige Art. Vi skiltes saa
ad uden at jeg sporede, at K. ikke kunde lide, at jeg havde
paataget mig den lille Rolle. I ethvert Tilfælde maatte han ogsaa
kende mit Forhold til sig tilstrækkeligt til at vide, at jeg ikke
kunde falde paa, ved at kopiere ham at bringe ham under en komisk
Belysning. Dertil havde jeg for megen Hengivenhed for ham, og ---
dertil var jeg meget for daarlig Skuespiller.

\secnum{24.}

I Begyndelsen af Fyrretyverne gjorde K. en Rejse over til
Vestkysten af Jylland, til hans Faders Fødeegn. Han fortalte mig
et lille karakteristisk Træk fra denne Rejse. Han var i sin
Fødeby, for hvilken hans afdøde Fader havde stiftet et større
Legat til en Skole, bleven modtagen med stor Ærbødighed og især
havde Skolelæreren paa Embedsvegne været stærkt i Aktivitet. „Jeg
var virkelig bange for, at man skulde rejse Æresporte for mig“,
sagde han. Da han skulde rejse, kørte han forbi Skolen. Der stod
Skolelæreren opstillet med alle Skolebørnene, der skulde afsynge
en af Læreren forfattet Sang til S. K.’s Ære. Læreren, der skulde
lede Sangen, stod med en Afskrift i Haanden og vilde netop give
Tegn til at begynde da K.’s Vogn holdt ved Siden af ham. S. K.
bøjede sig, med sit venligste Smil ned imod Skolelæreren, tog
Afskriften ud af Haanden paa ham, som for at læse den igennem, og
gav med det Samme Kusken et Vink om at køre til. Derved gik hele
Arrangementet i Stykker. Degnen, der ikke kunde sit Digt udenad,
kunde ikke sætte Sangen i Gang, Børnene stode tavse
% --- p. 355 ---
\opage{355}og forbavsede, og S. K. rullede hen ad Landevejen,
nikkende og hilsende, og i sit Hjerte morende sig over
Skolelærerens Skuffelse.

\secnum{25.}

Den Dag, da Licentiat Hagen (senere teologisk Professor) havde
disputeret, talte jeg med K. om hans Disputats („Om
Ægteskabet“), mod hvilken jeg havde opponeret. K. var kun lidet
tilfreds med den; han fandt hele Behandlingen overfladisk, og han
hentydede til, hvad der gives til en fyldigere Opfattelse af
Pseudonymerne. Ogsaa ved denne Lejlighed fremhævede han Tyskernes
grundigere Behandling af saadanne Opgaver, og sagde med Hensyn til
Hagens Afhandling: „man skulde dog sørge for, at vi kunne tage os
lidt anstændigt ud, naar saadanne Bøger blive læste af
Tyskerne.“

\secnum{26.}

Medens jeg i Sommeren 1846 opholdt mig i Berlin kom S. K. dertil.
Det var en af de smaa Rejser, han undertiden gjorde naar han ret
vilde arbejde intensivt, og søgte forandrede Omgivelser for at
sætte sig i Stemning. Jeg traf ham ganske uventet paa
Restavrationen, hvor jeg spiste, og hvorhen han nærmest var gaaet
for at se Landsmænd, da denne Restavration under hans første
Ophold i Berlin var de Danskes Samlingssted. Som Restavration
satte han ikke megen Pris paa den; den begrænsede sig, ytrede han,
til „det laveste Begreb af Afforing“. Han kom mig imøde med megen
Venlighed, indbød mig til at spise sammen med sig den næste Dag
paa Hotellet og gjorde lange Spadsereture med mig i det Par Dage,
han blev i Berlin. Jeg var meget optagen af alt det Nye, jeg i
Berlin kom i Berøring med og var utrættelig aktiv. K. morede sig
over denne Aktivitet, der ogsaa strakte sig til, at jeg besøgte
Møderne i en demokratisk Arbejderforening. Vi talte meget sammen
% --- p. 356 ---
\opage{356}om de Berlinske Forhold, om Forholdene hjemme, om min
Disputats, m. m. Paa vore Spadsereture lagde jeg Mærke til, at der
var enkelte Punkter, som K. havde faaet en Forkærlighed for, og
som han gerne opsøgte. Saaledes var der i Thiergarten et smukt
Parti med en Park omgiven af Blomsteranlæg. Den blev ham Maalet
for hans Spadseretur ligesom Blomsterpartierne ved Indgangen til
Frederiksberg Have havde været det i København. Maaske vakte den
Erindringer hos ham om denne og hvad der knyttede sig til den. ---
Ved at besøge ham i Hotellet lagde jeg Mærke til, med hvilken
sindrig Beregning Alt i hans Værelse var indrettet til at kunne
arbejde i Stemning. Belysningen, Forbindelsen mellem Værelserne,
Ordningen af Møblerne, Alt var ordnet efter en bestemt Plan. K.
stillede meget store Fordringer til Betjeningen, men var ogsaa
meget gavmild med sine Douceurer. Vor lille Diner til 2 Personer
serveredes i hans Værelse. Vi drak en Vin, han havde Forkærlighed
for paa Grund af Navnet: Liebfrauenmilch. Dette Navn havde givet
ham Indtrykket af noget Mildt og Let, og han drak Vinen i god Tro
til dette Indtryk uden at mærke, at det netop er en af de hedere
Rhinskvine. --- Det var mig velgørende at træffe sammen med K. paa
det fremmede Sted. Han bestyrkede mig i, hvad jeg dengang atter
begyndte at tænke paa, at vende tilbage til Danmark og der forsøge
paa at bryde mig en Bane. --- Fra Berlin sendte jeg engang til K.,
hvem Berlinernes Ueberschwänglichkeit altid morede, gennem vor
Kusine, Fru T. et Avertissement, jeg havde fundet i en Berliner
Adresseavis. Det lød, under Overskriften: Anfrage, saaledes:
„Sollte es nicht zweckmässig sein das auch die Klempner sich
assoziirten, um mit dem Zeitgeiste Schritt zu halten?“ --- Det
morede Kierkegaard, der genkendte sine Pappenheimere.

% --- p. 357 ---
% p. 357 opens a new section: the number stands at the head of the page.
\secnum{27.}

\opage{357}Efter at jeg var kommen hjem fra min første italienske
Rejse (i Slutningen af 1847) kom K. engang hen til mig,
da jeg paa Østergade stod og betragtede nogle Gibsafstøbninger
af Antiker, der stode i Vinduet hos Gibseren
Barsugli. Han ytrede, da vi kom til at tale om Billedhuggerkunsten,
at han egentlig ikke havde havt Øje for
den, men at den nu begyndte at tiltrække ham ved
sin Ro. Jeg tror overhovedet ikke, at K.’s Interesse og
Sans for den bildende Kunst var stærkt udviklet; hans
Interesse for det Skønne søgte aldeles overvejende sin
Næring i Poesien.

\secnum{28.}

Kort efter at Oprøret var udbrudt i 1848, talte jeg
med K. om den Plan, jeg dengang havde, at gaa med
som Frivillig. Jeg fremhævede, foruden den almindelige
Bevæggrund, der laa i Forholdene, det, at jeg personlig
trængte til at gøre Erfaring, navnlig Selverfaring. K.
gik ind derpaa, at denne Virkelighedens Lærevej for
Mange var nødvendig, men antydede, at han gennem
Fantasien og Refleksionen ideelt optog et endnu rigere
Indhold. Det forholdt sig ogsaa øjensynligt saaledes, at
% "Motiv" letterspaced: inter-letter gaps 13/9/10 px at 400 dpi.
han kunde nøjes med et \emph{Motiv} fra Virkeligheden, som
hans Refleksion og hans Fantasi forarbejdede, idet de
lededes af den Erfaring, han havde samlet gennem sit
eget Sjælelivs uophørlige Arbejden.

\secnum{29.}

Nogle Aar senere talte jeg med K. om et lignende
Æmne. Jeg følte mig dengang utilfreds ved det, ikke
at have en bestemt Virksomhed, og de bestemte Pligter,
der følge med en saadan; den fuldstændige Frihed, jeg
nød i mit privatiserende Studieliv, begyndte at blive mig
til Byrde. Jeg udtalte det i en Samtale til K. og sagde:
% „jeg: lower-case j as printed (checked at 500 dpi).
„jeg savner i den Grad en bestemt praktisk Virksomhed,
% --- p. 358 ---
\opage{358}at jeg, naar ingen anden findes, kunde fristes til at nedsætte
mig som Værtshusholder.“ — K. smilede og sagde:
„saa skal De faa min Søgning.“ Han fortalte mig i den
Anledning, hvorledes han nogle Aar tidligere havde faaet
et Besøg af en tysk Lærd, der rimeligvis af En eller
Anden var sendt hen til ham for at se paa ham som en
af vore Mærkværdigheder. Han havde modtaget Tyskeren
meget høfligt, men forsikret ham om, at der maatte være
en Misforstaaelse Skyld i Besøget. „Min Broder, Doktoren,“
havde han sagt, „er en overordentlig lærd Mand,
som det vist vilde interessere Dem at lære at kende,
men jeg er Ølhandler.“

\secnum{30.}

Dr. P. K. holdt engang, jeg husker ikke nøje Tiden,
% "Halvtredserne,": in the PDF's 1-bit mask the mark is a bare round point with no
% tail, and the second and third readers read a point. The text layer reads a comma,
% and IA's processed JP2 of this page (full-tone, 500 ppi) shows the comma's tail
% plainly: a COMMA prints. The tail was lost in thresholding the PDF image.
men det maa have været i Halvtredserne, under et tilfældigt
Ophold hernede, en Række Forelæsninger paa
Universitetet. De holdtes i Solennitetssalen for et talrigt
og meget blandet Avditorium. Søren K. forholdt sig
temmelig ironisk til disse Forelæsninger og det var med
en vis Malice at han fortalte mig et Brudstykke af en
% "Universitetet,": a sort stands here -- the line ends 17 px short of the measure, the width of a
% line-end comma (cf. "kende," above), and a displaced dot of ink lies where a comma's tail ends; the
% head did not print. Comma or point cannot be told apart even on IA's full-tone JP2 (UNDECIDABLE);
% the comma is read, since a point before lower-case "ligesom" would be a claimed defect without proof.
Samtale, han en Aften havde hørt udenfor Universitetet,
ligesom Forelæsningen begyndte. Der kom en Mængde
Herrer og endnu flere Damer strømmende op ad Universitetstrappen.
Kusken paa en Vogn, der holdt udenfor blev
tiltalt af en forbigaaende Karl med Spørgsmaalet om,
hvad alle disse Folk skulde derinde? „Aa! de skal vel
til Dans!“ svarede Kusken til stor Henrykkelse for S. K.
Forelæsningerne vare paa Doktorens Maner, temmelig
forceret aandrige, og i profetisk-apokalyptisk Stil. Jeg
kunde fortælle S. K. en Anekdote, Forelæsningerne vedkommende,
om vor gamle Onkel paa Købmagergade.
Han hørte regelmæssig Forelæsningerne, skønt de laa
meget over hans Horisont og skønt han var stærkt
tunghør. Han følte sig i høj Grad stolt over at se sin
„Fætters Søn paa Gammeltorv“ læse for saa talrig en
% --- p. 359 ---
\opage{359}Skare. En Aften, da han var kommen hjem fra Forelæsningerne,
gik han op og ned ad Gulvet i sin Stue,
uafbrudt gentagende: „nej! det er dog en mageløs Mand,
den Peter Kierkegaard! saaledes som han kan tale, saaledes
som han kan tale.“ Hans Datter tillod sig det
Spørgsmaal: „Kunde Du nu nok følge hans Foredrag?
% "en" letterspaced both times: gap 11 px and 10 px (two-letter words).
kunde Du rigtig høre hvad han sagde?“ — „Ikke \emph{en}
Ord! ikke \emph{en} Ord!“ sagde den gamle Jyde troskyldig,
men hans Begejstring over den mageløse Taler var lige stor.

\secnum{31.}

En kuriøs Situation med denne gamle Onkel fortalte
S. K. mig engang. Den gamle Onkel havde en Passion
for at gøre Vers, lige gyselige med Hensyn til Indhold
og Form. En Dag kom han op til S. K. og efter en
lille Indledning kom han frem med en Bunke af sine
Vers, som han bad Neveu’en at læse højt. De fik saa
Plads ved Siden af hinanden i Sofaen. Den Gamle sad,
lænet tilbage, og satte sine Briller paa for at kunne
følge med under Oplæsningen, nærmest vistnok for at
passe paa, at Intet blev sprunget over. Søren K. sad,
noget forover bøjet, med Papirerne, og valgte vistnok
denne Stilling for at den Gamle ikke skulde iagttage et
forræderisk Smil paa hans Ansigt. Han læste Digtene
op fra Ende til anden med hævet Stemme og med stor
Pathos. Den Gamle var aldeles henrykt, Taarerne løb
ham ned ad Kinderne af Rørelse over de skønne Vers,
og han forlod S. K. med de varmeste Taksigelser. —
Gamle Onkels Vers har jeg ofte hørt; de vare aldeles
mærkværdige, og havde navnlig den vidunderlige Egenskab,
med meget smaa Variationer at kunne bruges, eller ialfald
at blive brugte, ved de mest heterogene Anledninger.
Saaledes havde den Gamle engang, i Anledning af en
Datterdatters Forlovelse forfattet et Vers. Det blev saa
ikke blot benyttet ved en anden Forlovelse, hvor de Paagældendes
personlige Forhold dog vare meget forskelligt;
% --- p. 360 ---
\opage{360}men det skulde ogsaa med en ganske ubetydelig Tillæmpning,
have været benyttet i Anledning af min Ansættelse
ved Universitetet. At det ikke blev det, hidrørte kun
fra, at den Gamle ikke havde faaet det omordnet den
Aften, Udnævnelsen, lidt tidligere end vi havde ventet,
meddeltes af Berlingske, og vi netop var til Selskab hos
den Gamle. — Undertiden bleve Versene sungne af Selskabet
hos ham efter hans indstændige Anmodning. Der
valgtes saa en Melodi, der ved en Prokrustesmetode
nogenlunde kunde anvendes paa dem; men det var grundkomisk
at høre, hvorledes snart et stort Antal Stavelser
maatte ligesom tages i en Mundfuld, snart en enkelt
Stavelse trækkes ud efter den Holbergske Forskrift.
Men for den Gamles Ører klang det altid som de dejligste
Harmonier, hans Ansigt straalede af Henrykkelse
derover.

\secnum{32.}

Gamle Kierkegaard brugte undertiden i Tale og i
Skrift, dog især i Skrift, besynderlige Ordstillinger. Saaledes
plejede han, naar han i Brevene til sin Fødeby
% One pair of quotation marks only, opening at „min and closing after "S. K.",
% as printed.
omtalte S. K. og hans Broder, altid at skrive: „min
Fætters afdøde Søn Peter K., eller: min Fætters afdøde
Søn S. K.“ (i Stedet for: min afdøde Fætters Søn). Det
tog sig især komisk ud i én Sammenhæng: da han meddelte,
at hans Fætters afdøde Søn, Peter Kierkegaard,
havde taget sig en ny Kone. Jeg fortalte engang S. K.
om denne Udtryksmaade hos Onkelen. Den morede ham,
men for sit eget Vedkommende lagde han en for den
Gamle ubevidst Betydning (den Gamle var jo som Digter
inspireret!) ned deri: at han i Virkeligheden var en Afdød.
— Han kom ved samme Lejlighed til at tale om
Onkelens høje Alder. Jeg sagde, at i en vis Forstand
kunde man ikke kalde ham gammel, naar man nemlig
maalte Livets Længde efter dets Indhold; da den Gamle
i Virkeligheden kun „levede“ den mindste Del af den
% --- p. 361 ---
\opage{361}Tid han eksisterede, men mest blot vegeterede og sov.
Jeg føjede spøgende til, at egentlig var S. K. den ældste
Mand, jeg havde kendt. Han smilede og gjorde ingen
Indvending mod denne Regnemaade.

\secnum{33.}

Engang fortalte K. mig — det falder mig ind ved
at omtale hans Alder —, at han som ung i mange Aar
havde gaaet i den faste Overbevisning, at han skulde
dø naar han blev 33 Aar gammel. (Var det Jesu Alder,
der ogsaa skulde være Norm for Jesu Efterfølger?) —
Denne Tro var saa fast hos ham, at han, da han passerede
denne Alder, endog lod se efter i Kirkebogen for
at forvisse sig om, at den virkelig var overskreden; saa
vanskeligt var det for ham at tro dette.

\secnum{34.}

Da jeg første Gang holdt Forelæsninger ved Universitetet,
i Vinterhalvaaret 1849—50, talte S. K. hyppigere
med mig derom, om Forelæsningernes Genstand (den
græske Filosofi) og om Planen for Behandlingen. Han
sagde mig, at han engang vilde komme hen og høre
dem, men at han ikke vilde angive nogen bestemt Time.
% "dén": the accent over e is damaged (prints with an extra stroke) but is an acute.
Han kom dog ikke, men ved dén bestandige Udsigt til,
mulig at skulle have ham til Tilhører, gav han mig Impuls
til en omhyggelig Behandling af min Genstand, og
jeg tror nok, at det er det, han har tilsigtet. Han har
ogsaa vidst, at jeg let, naar han virkelig var tilstede,
vilde føle mig generet ved hans Nærværelse, og det kan
have været ham en Grund til at udeblive.

\secnum{35.}

I Anledning af mine Forelæsninger talte han engang
med mig om den aristoteleske Metafysiks Bestemmelse
% μεσότης: set in Greek type, smaller than the body.
af Dyden som μεσότης. Han havde rigtig indsét, at denne
Bestemmelse ikke gjaldt Dyden i Almindelighed, men
% --- p. 362 ---
\opage{362}kun den etiske Dyd, — Noget, man dengang ingenlunde
var paa det Rene med. Han fortalte mig, at han i Anledning
af dette Punkt og nogle andre i den Nikomachiske
Etik havde søgt Oplysninger hos Madvig og Sibbern.
Madvig havde sagt ham, at det var saa længe siden, at
han havde læst dette Skrift, at han ikke mere havde det
præsent. „Og Madvig er ellers et Jærn“, føjede S. K.
til. Sibbern derimod havde straks kunnet give ham de
søgte Oplysninger. S. satte S. K. megen Pris paa, om
han end ikke var blind for hans Svagheder. Blandt
disse fremhævede han engang Sibberns fuldstændige
Mangel paa Ironi, og — i psykologisk Henseende —
hans Mangel paa Blik for de maskerede Lidenskaber,
for den Fordobling, ved hvilken den ene Lidenskab
tager den andens Form. Derved var, efter hans Mening,
Sibbern tidt blevet skuffet paa en Tid, da mange, navnlig
Damer, henvendte sig til ham for at konsulere ham som
en Art psykologisk Sjælesørger.

\secnum{36.}

Poul Møller kom S. K. hyppig til at tale om, altid
med den inderligste Hengivenhed. Det var Poul Møllers
Personlighed, langt mere end hans Skrifter, der havde
gjort Indtryk paa ham, og han beklagede, at der snart
vilde komme en Tid, da man, naar den bevarede Erindring
om hans Personlighed var forsvunden, og Dommen
om ham kun støttede sig paa hans Skrifter, ikke mere
vilde kunne forstaa hans Betydning. Han fortalte mig
engang et morsomt lille Træk af P. Møller. Han skulde
ex officio opponere ved en Disputats og havde optegnet
sine Bemærkninger paa forskellige løse Blade, der vare
lagte ind i Disputatsen. Hver enkelt Indvending indledte
han med et „graviter vituperandum est“, men saasnart
Præses havde givet et Svar paa Indvendingen, sagde han
meget godmodig: Concedo, og gik over til den næste
Indvending. Efter en temmelig kort Opposition sluttede
% --- p. 363 ---
\opage{363}han med en levende Beklagelse over, at hans Tid ikke
tillod ham at fortsætte denne interessante Samtale, og
idet han saa fjærnede sig og gik forbi S. K., der stod
blandt Tilhørerne, sagde han halvhøjt til ham: „Skal vi
saa gaa ned til Pleisch?“ — det var den Conditor, hvor
han plejede at søge. — Under Oppositionen var det hændet
ham, at alle hans løse Papirer paa én Gang vare faldne
ud af Bogen og flagrede omkring paa Gulvet, og det
havde ikke lidet bidraget til at forhøje Avditoriets Stemning
at se den store Skikkelse kravle omkring for at
opsamle de spredte Blade.

\secnum{37.}

En Mand, med hvem jeg i de Aar jævnlig saa’ K.
spadsere, var Professor Kolderup-Rosenvinge. Det var
nærmest en fælles Interesse for det Æstetiske, der bragte
dem i Berøring. K.-R. var en dannet Mand, der beskæftigede
sig med de sydevropæiske Literaturer og havde
oversat et Stykke af Calderon. Forresten var han dengang
allerede temmelig sløvet og i mange Henseender
yderlig bornert. Da jeg engang for S. K. udtrykte min
Forundring over, at han kunde finde Interesse i at tale
med Rosenvinge, fremhævede han hans almindelige Dannelse.
Han satte overhovedet megen Pris paa Folk af
den ældre Generation, der havde bevaret den ældre Tids
humane Interesser og den Finhed i Væsen, som en yngre
Generation saa meget savnede.

\secnum{38.}

Afdøde Pastor Spang havde K. en Tid lang (i Aarene
1844—45?) omgaaedes meget, og havde næsten daglig
afhentet Spang i hans Hjem til en længere Aftenvandring.
(Adskillige interessante Breve til Spang fra S. K. findes
hos Spangs Datter, Fru Rump). Efter at Spang var død,
fortalte K. mig engang, hvorledes hans Enke havde været
overvældet af Sorg, og hvorledes han havde ved sin
% --- p. 364 ---
\opage{364}Samtale bidraget til, atter at bringe hendes Sind i Ligevægt.
Han brugte ved den Lejlighed det Udtryk: „Jeg
véd, hvorledes man skal sørge, og hvorledes man skal
trøste i Sorgen.“ Og tilvisse: han forstod det som Faa.
Han trøstede, ikke ved at tildække Sorgen, men ved
først at bringe den ret til Bevidsthed, at klare den helt
% "Pligt" letterspaced (gaps 11/10/11/10 px); the second "Pligt" is not.
ud, og saa at minde om at det var en \emph{Pligt} at sørge,
men netop derfor ogsaa en Pligt, ikke at lade sig knuge
ned af Sorgen, men i Sorgen at bevare Kraften til sin
Gerning, ja i den at finde Incitamentet til endnu fuldstændigere
at gøre denne Gerning. — Jeg har selv i mit
Forhold til S. K. mere end én Gang gjort den Erfaring,
hvorledes han forstod at løfte En naar man var nedbøjet,
% "det uden": a comma-shaped blob of ink lies wholly below the line after "det", but no comma
% sort is set: the gap to "uden" is an ordinary word space on this line (a comma takes its own
% width). Confirmed on IA's full-tone JP2: no comma prints. Stray ink; not transcribed.
trøste naar man sørgede, og det uden at man havde
nødig at tale om hvad der kuede eller gjorde bedrøvet.
Saaledes mindes jeg engang i Efteraaret 1850, da jeg,
fordi jeg ikke i rette Tid havde kunnet finde en Lejlighed,
en kort Tid maatte bo paa en uhyggelig Gæstgivergaard,
% PRINTER'S ERROR: "Iudtrykket" printed with a turned n (for "Indtrykket").
og under Iudtrykket af Omgivelserne de mismodige Tanker
fik Overhaand, at jeg da en Aften, medens jeg gik forstemt
paa Gaden, mødte K., der indledte en Samtale
med mig. Uden at jeg behøvede at tale et Ord om mig
selv, saa’ han med sit skarpe Blik, at jeg trængte til
at bringes ud af en kuende Stemning, og han forstod
ved sin Samtale, uden at den et Øjeblik direkte syntes
rettet derpaa, saaledes at frigøre mit Sind, at jeg forlod
ham glad og frejdig, og for lang Tid var befriet for
Tungsindets Magt.

\secnum{39.}

Meget af det som S. K. har udført i sine Skrifter,
mindes jeg, at han i Samtale med mig har berørt medens
hans Tanke arbejdede med det, og det fik netop derved,
at det endnu var i sin Vorden, et Liv og en Tiltrækning
endnu større end den, det kunde have som fuldt udformet.
Saaledes mindes jeg, at han engang, medens han fulgtes
% --- p. 365 ---
\opage{365}med Christens og mig, udførte det Tema, han saa tidt
% "hele" letterspaced (gaps 10/10/11 px) -- the brief's calibration word.
har behandlet i sine opbyggelige Skrifter, at „\emph{hele} Livet
er Prøvelsens Tid.“ Han udførte det i den Form, at
der for den Troende var et „Omsider“, der syntes at
være saa nær, at det kunde gribes, men bestandig veg
tilbage, indtil Livets Slutning endelig førte til at gribe
det. Først i Døden kunde der siges: „Omsider!“ —
Undertiden udtrykte S. K. i sine Samtaler sin Tanke i
en mere tilspidset, mere paradoks Form end i sine Skrifter.
Det morede ham at formulere den saaledes, at den dannede
Modsætningen til et almindelig Vedtaget. I Modsætning
til Sætningen: Erfaringen gør klog, plejede han saaledes
at opstille den Teori: „Erfaringen gør gal!“ og han
kunde udføre Beviset med megen Lune ved at vise hen
til alle den konkrete Erfarings utallige Modsigelser.

\secnum{40.}

Den forskellige Retning i vort aandelige Liv og vort
grundforskellige Forhold til det Kristelige kom af og til
% PRINTER'S ERROR: "sammen, Engang" -- a comma where the sentence ends
% (capital E follows); checked at 500 dpi, the tail is clear.
bestemt tilsyne naar vi talte sammen, Engang, vistnok
i 1851, spurgte han mig, hvad jeg for Øjeblikket beskæftigede
mig med. Jeg svarede ham, at jeg læste det
N. T. Det syntes at behage ham, men det behagede
ham ikke da jeg tilføjede, at jeg nærmest læste det med
det Maal for Øje, i det Overleverede at finde det primitivt
Kristelige og at følge den sukcessive Udvikling af de
kristelige dogmatiske Forestillinger. For ham var en
saadan Undersøgelse noget næsten Forargeligt. Skønt
han i sine Skrifter saa lidt drager Dogmerne ind og overalt
holder sig til det Centrale i Kristendommen, kunde
han dog ikke stille sig i et saadant Forhold til Skriften,
at den blev Genstand for en kritisk Undersøgelse. Han
kunde det ikke, dels fordi Skriften for ham stod som en
Enhed, som et Enhedsudtryk for det Kristelige, i hvilket
han ingen Sondring vilde bringe ind; dels fordi den, ved
% --- p. 366 ---
\opage{366}at gøres til Genstand for en videnskabelig Undersøgelse,
blev Vidensobjekt i Stedet for Trosobjekt.

\secnum{41.}

% PRINTER'S ERROR: "Forboldet" printed with b for h (for "Forholdet").
Min Afvigelse fra S. K. i Forboldet til det Kristelige
skjulte jeg ikke for ham; men vi disputerede ikke
derom. Der var saa overordentlig Meget, i hvilket vi
kunde sympatisere, og jeg foretrak at lade mig belære
af ham om Meget, der var mig af stor Betydning,
for at føre en Diskussion, i hvilken han ved sin Overlegenhed
let vilde have kunnet overvælde mig, men ikke
let overbevist mig. Han forholdt sig ogsaa i vore Samtaler
væsentlig meddelende til mig, og jeg saa’ deri et
Bevis paa, at han nærede Velvillie imod mig. Han havde
to Maader at samtale paa: den ene som væsentlig meddelende,
vækkende, opmuntrende; den anden som spørgende,
ironiserende, og gennem sin Dialektik konfunderende.
Den sidste har han aldrig anvendt imod mig.
En enkelt Gang kunde han, naar jeg blev ivrig, kaste et
spøgende, drillende Ord ind. Saaledes talte jeg engang
meget ivrigt om, at ingen positiv Religion kunde være
tolerant, netop fordi den, ved sit Krav paa at være
aabenbaret Religion maatte gøre Fordring paa at være
den eneste sande, og derfor maatte sætte de andre som
usande. Et alment Religiøst, en „Religion overhovedet“
maatte derfor, fra den positive Religions Synspunkt, være
en Uting. Idet jeg med Ivrighed udviklede dette, kom
jeg til at gentage Udtrykket: en Religion overhovedet,
og gjorde den Sætning til min, at en Religion (d. v. s.
% "positiv" letterspaced (gaps 13/15/11/11/9/10 px).
\emph{positiv} Religion) overhovedet var en Uting. — „Ja, og
en Natpotte over Hovedet“ sagde K., og dæmpede derved
min Ivrighed. Det er imidlertid det eneste Drilleri af
den Art, jeg mindes af ham. — Naar han aldrig prøvede
paa, direkte at imødegaa mine afvigende Anskuelser,
kan jeg, efter den hele venlige Deltagelse, han altid viste
mig, og som jeg taknemmelig erindrer, kun søge Grunden
% --- p. 367 ---
% p. 367 opens mid-sentence (not mid-word), flush left: no new paragraph.
\opage{367}til det deri, at han vidste, at jeg ærligt og alvorligt beskæftigede
mig med den Sag, der for ham var den højeste, og at jeg var saa fortrolig med
hvad han havde skrevet derom og med hans hele Tankegang, at jeg havde
tilstrækkelige Præmisser til en Konklusion, og ikke vilde blive bragt den
nærmere ved en Anden, men ialfald kun ved mig selv.

\secnum{42.}

% "ham" (line 3 of section): h damaged at the head, reading certain.
Om den almindelige Virkning af S. K.’s Skrifter sagde jeg engang et Ord til
ham, der syntes at gøre noget Indtryk paa ham. Han omtalte, hvorledes han som
Forfatter havde havt den ejendommelige Skæbne, at udgive det ene Skrift efter
det andet uden at fremkalde en egentlig Kritik, angribende eller rosende. Han
fortalte saa, at hans Skrifter dog fra Begyndelsen af havde fundet en Kreds af
Læsere, som det syntes en fast Kreds, da Antallet af de Eksemplarer, der
straks afsattes, omtrent var det samme (c. 70 for de Skrifter, der fulgte
efter Enten-Eller). Jeg udtalte min Overbevisning om, at overmaade Mange
% a raised stray dot (a spacing sort) stands between "hans" and "Skrifter"; not transcribed.
læste hans Skrifter med virkelig Alvor, og at vor Presses Tavshed, som jeg
nærmest forklarede ved Kritikernes Følelse af at Skrifterne
% "laa'" as printed, with apostrophe (past tense, cf. "saa'" p. 371).
laa’ over deres Kritiks Kræfter, ingenlunde var et Tegn paa, at han stødte paa
Ligegyldighed hos et større Publikum, paa hvilket jeg netop troede, at han
allerede dengang indvirkede efter en stor Maalestok. Men jeg føjede til, at
Virkningen med Hensyn til Forholdet til Kristendommen efter min Mening havde
været ligesaa meget negativ som positiv. Alle, paa hvem Skrifterne i den ene
eller den anden Retning havde havt en Indvirkning, vare vistnok enige om, at
Kristendommen kun i hans Opfattelse havde en virkelig Ejendommelighed og
indgød Ærbødighed. Men hans Bestemmelse af Kristendommen var en saadan, at den
netop maatte kunne bort%
% --- p. 368 ---
\opage{368}skyde Mange fra den, saaledes som det var Tilfældet med mig, og den
maatte kunne det netop ved den Adskillelse, den satte mellem Kristendommen og
Naturen og det konkrete Menneskeliv. --- Vi vare, da jeg sagde disse Ord til
K., netop ifærd med at tage Afsked fra hinanden; han blev alvorlig og gav mig
tavs Haanden, og skiltes fra mig, men med et venligt Udtryk i sit Ansigt.

\secnum{43.}

Uforbeholdent udtalte S. K. sig om den Betydning, hans Skrifter kunde have
for visse bestemte Personer. Hans gamle Onkel, Grosserer M. Kierkegaard, havde
en Søn, et Par Aar yngre end S. K., der var vanfør, lam i den hele ene Side,
og fuldstændig Krøbling; aandelig derimod vel begavet. Han læste med den
største Interesse Fætterens Skrifter, besøgte ham af og til i hans Hjem, og
hentede sig ved dette Besøg megen aandelig Løftelse. Om ham talte jeg engang
med S. K. og fortalte ham, hvor stort et Indtryk et af hans Arbejder:
Skriftetalen i „Opbyggelige Taler i forskellig Aand“, havde gjort paa ham. I
den tales der et Sted om et Menneske, der ved sin legemlige Svaghed er
forhindret fra at udøve en Gerning i det Ydre, og der skildres smukt og
opløftende, hvorledes ogsaa et saadant Menneske har den
% "almen-/menneskelige" broken at a line end; joined as one word -- the hyphen may belong to the compound.
almenmenneskelige etiske Opgave ubeskaaren, og hvilken særlig Form
% "ham.": one point. A small low blob follows it (about 40% of a point, 2 px below the
% baseline) and takes no set width -- the space to "S." is the line's word space plus the
% usual sentence increment. A speck, not a second point (third reading; IA's JP2 agrees).
% The text layer reads "ham...".
Livsopgaven faar for ham. S. K. sagde: „Ja! for ham er denne Bog en Lykke.“
Det var det ogsaa virkelig. Det havde Magten til at give den haardt Prøvede
Kraft til at overvinde den Tanke, at hans Tilværelse var unyttig og forspildt,
og at give ham Bevidstheden om, væsentlig at være lige med de af Naturen
lykkeligere Udrustede. Netop denne Bevidsthed forstod S. K. ogsaa ved sine
Samtaler at kalde tillive hos ham; derfor gik han saa styrket fra dem.

% --- p. 369 ---
\secnum{\opage{369}44.}

% emphasis: "nu er jeg færdig" letterspaced; at 400 dpi letter gaps 9--15 px
% ("færdig" 7,10,11,9,12), plain "Denne" on the same line 4--5 px.
En Samtale, jeg siden hyppigt har mindedes, havde jeg med S. K. en Tid før
hans „Indøvelse i Kristendom“ udkom. Vi gik sammen nede ved Søerne og han
opfordrede mig til, ikke for længe at opsætte det, at begynde en
Forfattervirksomhed. Han føjede spøgende til: „Nu er der jo en Plads ledig i
vor Literatur, \emph{nu er jeg færdig}.“ Denne Ytring, som jeg dengang gik let
hen over, kom sildigere levende frem i min Erindring da det sidste Afsnit af
S. K.’s Forfattervirksomhed laa afsluttet. Da han talte saaledes til mig, har
han øjensynlig endnu næret det Haab, at kunne undgaa den sidste Kamp med det
Bestaaende. Var den ikke bleven fremtvungen, havde hans Forfattervirksomhed
ogsaa været --- væsentlig endt, afsluttet, da hin Samtale fandt Sted. I
Virkeligheden vare alle Stadierne, der skulde gennemløbes, dengang
gennemløbne; hans Skrifter dannede et forunderligt afsluttet Hele. Men det
næste Stadium bar endnu en Mulighed til en Fordobling i sig, og denne
Mulighed blev til Virkelighed, da der gjordes et Forsøg paa at reducere hans
Opfattelse af Kristendommen til en ekscentrisk Overdrivelse, at sætte
Middelmaadigheden og Aandløsheden som det egentlig Normale.

\secnum{45.}

I den samme Samtale kom S. K. til at fremhæve en Side ved sin
Forfattervirksomhed, som han omtalte med den ham egne Objektivitet: det var
hans Skrifters Betydning for den danske Prosa. Han sagde: „Vor Literatur har
i dette Aarhundrede fremvist en næsten abnorm Rigdom i Poesiens Udvikling; vor
Prosaliteratur derimod har staaet tilbage; vi savnede en Prosa, der havde det
Kunstneriskes Præg. Dette Savn har jeg udfyldt, og den Betydning ville mine
Skrifter ogsaa bevare for vor Literatur.“ Heri tror jeg at han har Ret: hans
Prosa er Kunst, ikke uden Skygger, men som Helhed det klareste,
% --- p. 370 ---
\opage{370}bøjeligste, omfangsrigeste og fyldigste Udtryk for Tanken, og ikke
mindre for Følelsen og Lidenskaben, som vor Literatur kender. Den har hyppigt
mindet mig om Plato, den Forfatter, med hvis Prosa jeg nærmest vilde
sammenligne Kierkegaards, og som vistnok ogsaa har givet ham Forbilledet, som
han frit har efterdannet.

\secnum{46.}

I S. K.’s Skrifter har han paa flere Steder omtalt Grundtvig og
Grundtvigianismen, og saagodtsom overalt set den fra den komiske Side. Paa
lignende Maade udtalte han sig hyppig om den i sine Samtaler. Den skarpe
Kontrast, der er imellem hans Opfattelse af det Religiøse og den grundtvigske,
traadte bestandig frem, og den grundtvigske barnagtige Umiddelbarhed fik i
denne Kontrast en komisk Belysning. Kun om et Afsnit i Grundtvigs Liv har
S. K. engang talt til mig med Interesse og Anerkendelse: det var om den Tid,
da Gr. optraadte mod Clausen. Men hvad der i hans Optræden dengang tiltalte
K., var øjensynlig ikke det specifik Grundtvigske, allermindst den „mageløse
Opdagelse“, der altid forekom K. at være en mageløs Taabelighed og tilmed en
ikke original Taabelighed, men det, der danner Modsætningen til Rationalismen:
% emphasis: only "Livs" is letterspaced (gaps 10,10,11, and 11 before "r"); "rørelse" 3--6 px, as plain.
den religiøse \emph{Livs}rørelse. S. K.’s Fader, hvis religiøse Retning
nærmest var den gamle Pietismes, sympatiserede med Grundtvig og Lindberg,
vistnok fra dette Synspunkt, og S. K.’s Sympati havde upaatvivlelig samme
Motiv og samme Grænser som Faderens.

\secnum{47.}

Et lille Træk af S. K. falder mig ind ved at tænke paa Grundtvig.
Grundtvigianismens Patronesse, Enkedronning Caroline Amalie, satte ogsaa megen
Pris paa S. K., idetmindste før hans Angreb paa Mynster. S. K. var øjensynlig
ikke ufølsom for en Velvilje, der vistes
% --- p. 371 ---
\opage{371}ham fra en Dronning. En Dag gik jeg med ham paa Kærlighedsstien, da
han pludselig udbrød: „Naa, for Fanden! nu slipper jeg ikke.“ Jeg saa’
forundret paa ham, der ellers aldrig brugte Ord, der lignede en Ed, og saa’
frem ad Stien, hvor jeg i lang Afstand øjnede en rød Lakaj, og foran ham to
Damer. Den ene af disse var Enkedronningen, der, da vi mødte hende, meget
huldsalig standsede S. K. og talte med ham. Vi var, da han saa’ hende, saa
langt fra hende, at vi godt kunde være vendte om, uden at være blevne
bemærkede. At hans Udbrud ikke var aldeles oprigtig ment, fik jeg et Indtryk
af ved dette: „naa, for Fanden!“ det lignede ham ikke, og det kom noget
tvungent, som naar man ved en stærk Ytring vil skjule, at man er ret vel
tilfreds med at Noget skér. S. K.’s gamle, traditionelle Ærbødighed for
Kongemagten har i hans Forhold til Enkedronningen ikke været uden Betydning.

\secnum{48.}

Om Biskop Mynster havde jeg, i Foraaret 1852, en Samtale med S. K., som jeg
siden hyppigt maatte tænke paa, da K. optraadte imod ham. Jeg gik en
Eftermiddag paa Volden, paa den Spadseresti, der er oppe paa Brystværnet, og
gik der forbi Biskop Mynster. Et Øjeblik efter mødte jeg S. K., der gik nede
paa Volden og kaldte paa mig. Han havde set Mynster gaa foran sig, og derved
faldt Samtalen mellem os paa ham. K. omtalte først hans anselige og
ærværdige Ydre, og talte saa om hans kloge og fine Optræden, og hvorledes han
som Verdensmand var de Fleste af de Yngre overlegen, der en Tidlang havde
troet, med deres filosofiske Apparat at kunne behandle Mynster som en
underordnet Figur, men snart, ved personlig Berøring med ham, vare blevne
overvældede af ham og nu underordnede sig under ham. Han nævnede navnlig
% five spaced full points printed in place of a name.
. . . . . som den, der ved sin Kejtethed og Mangel paa Finhed gav Mynster en
saa høj Grad af
% --- p. 372 ---
\opage{372}Overlegenhed i det personlige Forhold. Men han betonede meget stærkt,
at Mynsters Force netop var hans Verdensklogskab, og at det var den, der
bestemte hans Optræden. Angaaende sit eget Forhold til Mynster talte K.
ironisk. Han fortalte mig med et Smil, at han engang havde anmodet Mynster om
Tilladelse til, naar der var Problemer, som han ikke selv kunde bringe til
Klarhed, da at søge Vejledning hos hans højærværdige Ekscellence! og at
Mynster velvillig havde givet ham denne Tilladelse. Alene det Smil, hvormed
han sagde mig dette, viste mig tydeligt, hvorlidet højt han i Virkeligheden
dengang satte Mynster, og den hele Maade, paa hvilken han karakteriserede
Mynster som Verdensmand, stemmede fuldkomment overens med hans senere
Udtalelser. Der var en Tid, da han havde sat Mynster meget højt, og dertil var
han væsentlig kommen gennem Pieteten mod sin Fader, der satte stor Pris paa
Mynster. Dette pietetsfulde Hensyn til Faderens Dom har jeg i flere Tilfælde
set, som bestemmende K.’s Opfattelse. Saaledes mindes jeg, at han engang med
megen Velvilje omtalte en Mand, der i ingen Henseende fortjente Agtelse eller
ved sin Person vilde kunne indgyde S. K. Interesse; men hans Fader nærede stor
Hengivenhed og Erkendtlighed mod denne Mand, der engang ved et klogt Raad havde
forhindret, at hans Pengeaffærer, der i den forvirrede Pengeperiode efter
Krigen mod England vare truede, kom i Forstyrrelse.

\secnum{49.}

Dengang S. K. begyndte sin Polemik mod det Bestaaende, og maaske allerede
nogen Tid tidligere, ophørte han at deltage i den kirkelige Gudstjeneste, ved
hvilken han tidligere havde været en meget stadig Gæst. Afdøde Dr. Frederik
Beck fortalte mig engang, at S. K. i den Tid gærne kom i Athenæum Søndag
Formiddag under Kirketiden; hvilket Beck, med den for ham ejendommelige Lyst
til spydige Bemærkninger, fortolkede
% --- p. 373 ---
\opage{373}saaledes, at S. K. vilde gøre Folk opmærksomme paa, at han ikke
besøgte Kirken. --- Engang, mange Aar tidligere, havde S. K. spurgt mig, om
jeg gik i Kirke. Jeg benægtede det, og sagde, at jeg ikke gjorde det, dels
fordi jeg følte mig fremmed der, dels fordi Gudstjenesten i de københavnske
Kirker foregik under saamange Forstyrrelser, at jeg i intet Tilfælde kunde
bevare en andægtig Stemning, men snarere maatte forarges. Han sagde da, at
% emphasis: "Pligt" letter gaps 10,9,9,10 px; "skulde" 10,9,11,11,13 px; plain words on the lines 2--6 px.
det var en \emph{Pligt}, ikke at lade sig forstyrre ved hvad der foregik
omkring En; man \emph{skulde} opbygges naar man var i Kirke.

\secnum{50.}

Sidste Gang talte jeg med S. K. i Sommeren 1855, efterat jeg usædvanlig længe
ikke havde sét ham. Det var en Aften, da jeg gik fra Højbroplads op imod
Vimmelskaftet. K. kom hen til mig og indledte en Samtale. Efterat vi et
Øjeblik havde talt om min kortvarige Deltagelse i det politiske Liv (jeg
% "1854---55": the dash ends the line, "55" begins the next.
havde i Vinteren 1854---55 været Medlem af Rigsdagen), kom vi snart til at
tale om hans Polemik mod det kirkeligt Bestaaende. Jeg takkede ham varmt for
hans Optræden, og sagde ham, hvor Mange jeg havde fundet, der inderlig
sympatiserede med ham. Han talte om den Situation, han havde fremkaldt, med den
største Klarhed og Ro, og det forbavsede mig, at han under den voldsomme Kamp,
der greb saa dybt ind i hans Liv og gjorde Krav paa hans sidste Kræfter, kunde
bevare ikke blot den vante Sindsro og Frejdighed, men endnu kunde bevare
Spøgen. Da vi nærmede os til hans Hjem --- han boede dengang i Klædeboderne ---
sagde han i en spøgende Tone: „Nu gaar jeg hjem og gaar til Ro, og saa siger
jeg til Verden ligesom hin Kammerraad.“ Jeg spurgte ham, hvad denne Kammerraad
da havde sagt. Han fortalte mig saa, at der engang havde boet en Særling der i
Gaden, der var Kammerraad af Titel, og som regelmæssig hver Aften
% --- p. 374 ---
\opage{374}havde ført følgende Samtale med Vægteren. Naar Vægteren havde raabt
10, aabnede Kammerraaden Vinduet, og spurgte: „Vægter, hvad er saa Klokken?“
% printer's error: the opening quote before "Godt!" is two separate commas (,,), not the „ sort used elsewhere; transcribed „.
% printer's error: this quotation („Godt! ...) is never closed -- no closing “ before or after the parenthesis.
--- „Klokken er 10, Hr. Kammerraad.“ --- „Godt! saa tror jeg, at jeg vil gaa
til Sengs. Hvis nu Nogen spørger efter mig, Vægter, saa kan De bede dem at ---
% "Götz": only ONE dot prints over the o -- round, 7x8 px, foot 5 px above the letter, over its left third,
% i.e. the size and position of the left dot of a diaeresis (cf. "ä" of "Ueberschwänglichkeit" p. 356),
% not of this fount's accents (the acute of "sét" p. 373 is a 13x13 px wedge). The right dot did not
% print. IA's full-tone JP2 cannot decide between ö and ò; ö is read, a grave being a claimed defect
% without proof. The text layer reads "Gotz".
(Götz v. Berlichingens Svar.) --- Han fortalte denne Historie med sit gamle
Lune, og jeg sagde ham Godnat! uden at ane, at jeg den Aften skulde have talt
med ham for sidste Gang.

\begin{center}\small (Nedskrevet i Tiden fra 27. Dec. 1871 til 10. Januar 1872.)\end{center}

\begin{flushright}\textbf{H. Brøchner.}\end{flushright}

% End of article. Below the signature p. 374 carries only a short centred rule; the rest of the page is blank.
% p. 375 (PDF 385, no folio) opens the next article, „Om nyere spansk Literatur.“ (bold sans-serif head over a short rule) -- not transcribed.

\end{document}
